% Study guide for the structural-idealization / multi-cell torsion series. Build: `mtk guide guide.tex`
% (after `mtk episodes`). Every number here is computed and asserted in check_guide.py.
\documentclass[11pt]{article}
\usepackage[margin=0.9in]{geometry}
\usepackage{amsmath,amssymb}
\usepackage{mtkguide}
\usepackage[hidelinks]{hyperref}
\setlength{\parindent}{0pt}
\setlength{\parskip}{0.4em}
\newcommand{\mm}{\ \mathrm{mm}}
\newcommand{\Nmm}{\ \mathrm{N/mm}}

\begin{document}
{\LARGE\bfseries Booms, Shear Flow and Multi-Cell Torsion}\par{\large Practice Guide}\par
{\color{mtkmuted} AERSP 301, HW 5 topics \textperiodcentered{} a guide to work through with the video}\par\medskip

\textbf{How to use it.} The video pauses at each checkpoint. Work it here, on paper or a tablet, before you press
play. Checkpoints 2.1, 3.1 and 5.1 are answered by the video as soon as you resume. Checkpoints 4.1 and 6.1 are
fresh practice problems that the video never solves: they repeat the homework's steps with new numbers, so check
them against the full solutions on page~\pageref{solutions}. Stuck for five minutes? Read the hint on
page~\pageref{hints}.

\section*{Episode 2: Why Booms?}

\checkpoint[2.5cm]{2.1}{A web in bending}{%
A vertical web $300\mm$ deep and $10\mm$ thick is symmetric about the neutral axis of a beam in pure bending, and
carries direct stress. When the section is idealized, how much area does this web add to the boom at each of its
ends? Use $B_1 = \dfrac{t_D b}{6}\left(2 + \dfrac{\sigma_2}{\sigma_1}\right)$.}

\section*{Episode 3: Shear Flow With Booms}

\checkpoint[3.5cm]{3.1}{Finish the channel}{%
The channel has four booms of $300\mm^2$: booms 1 and 2 at $y = +200\mm$ (top flange), booms 3 and 4 at $y = -200\mm$ (bottom flange), and $S_y/I_{xx} = 10^{-4}\ \mathrm{N/mm^4}$. Starting
from $q = 0$ at the free edge, the flow in the top flange is $q_{12} = -6\Nmm$. Using
$q_2 - q_1 = -\dfrac{S_y}{I_{xx}}B_r y_r$, find $q_{23}$ (the web), $q_{34}$ (the bottom flange) and the flow just
past boom 4. What should the last one be, and why?}

\section*{Episode 4: Homework Problem 1}

\checkpoint[9cm]{4.1}{Practice box: booms, then the shear centre}{%
A box section is $600\mm$ wide and $240\mm$ deep (centreline dimensions). Top and bottom skins are $8\mm$ thick, the
left web $6\mm$, the right web $5\mm$. The left corners carry $50\times50\times6\mm$ angles and the right corners
$40\times30\times5\mm$ angles. One material throughout.
(a) Idealize it into four corner booms (skins and webs carry direct stress: use the ratio $+1$ for the skins and
$-1$ for the webs in bending). (b) Find $I_{xx}$. (c) Number the booms as in the video (1 top left, 2 top right, 3 bottom right, 4 bottom left), cut the top wall,
go $1\to2\to3\to4\to1$, and find the basic shear flow $q_b$ in each wall for an upward load $S_y$. (d) Use zero twist to find $q_{s,0}$, then find the shear centre's distance $\xi$ from the
left web.}

\section*{Episode 5: Torsion of Closed Sections}

\checkpoint[4cm]{5.1}{Twist rate of one cell}{%
A two-cell box carries shear flows $q_I$ and $q_{II}$, both taken counter-clockwise. Cell I has area $A_I$; the wall it shares with cell II
has $\delta_{12} = \int ds/t$, and all of cell I's walls together (shared wall included) have $\delta_I$. Write
$d\theta/dz$ for cell I in terms of $q_I$, $q_{II}$, $\delta_I$, $\delta_{12}$, $A_I$ and $G$.}

\section*{Episode 6: Homework Problem 2}

\checkpoint[9cm]{6.1}{Practice two-cell box: $T = 15$ kN m}{%
A two-cell box is $250\mm$ deep. Cell I (on the left, with the left web) is $500\mm$ wide and cell II is $600\mm$ wide
(rectangular cells, centreline dimensions; the spar between them is shared). Walls: left web $2.0\mm$, spar $2.5\mm$, right web $1.6\mm$, all with
$G = 27\,000\ \mathrm{N/mm^2}$; cell I's top and bottom skins $1.0\mm$ and cell II's top and bottom skins $0.8\mm$, all
with $G = 24\,000\ \mathrm{N/mm^2}$. It carries a counter-clockwise torque $T = 15\ \mathrm{kN\,m}$.
Using $G_{REF} = 27\,000\ \mathrm{N/mm^2}$, find (a) $t^*$ and $\delta$ for each wall, (b) $q_I$, $q_{II}$ and the
spar's shear flow, (c) the shear stress in each wall, $\tau = q/t$ with the real thickness, and (d) the rate of twist from
$G_{REF}\,d\theta/dz = \frac{1}{2A_I}(q_I\delta_I - q_{II}\delta_{spar})$, in degrees per metre.}

\clearpage
\section*{Hints}\label{hints}
\hintfor{2.1} In pure bending, what are the stresses at the top and bottom of a web that is symmetric about the
neutral axis? That gives $\sigma_2/\sigma_1$.
\hintfor{3.1} Boom 3 is below the axis: its $y_r$ is negative. A free edge has nothing to pass shear flow on to.
\hintfor{4.1} An angle's area is (leg + leg $-$ thickness) $\times$ thickness. Only the two webs have a basic flow;
the cut wall and the opposite skin come out zero. For $\xi$, the web forces must add up to $S_y$.
\hintfor{5.1} Go around cell I: which walls carry $q_I$, and which one carries $q_I - q_{II}$? Collect the $q_I$ terms.
\hintfor{6.1} Scale only the skins: $t^* = (G/G_{REF})t$. Set cell I's twist rate equal to cell II's, then use
$T = 2A_Iq_I + 2A_{II}q_{II}$ (with $T$ in N\,mm).

\clearpage
\section*{Solutions}\label{solutions}

\solutionfor{2.1}{A web in bending}
In pure bending the web's two ends sit equal distances above and below the neutral axis, so their stresses are equal
and opposite: $\sigma_2/\sigma_1 = -1$.
\[ B = \frac{10\times300}{6}\,(2 - 1) = 500\mm^2 \text{ to each boom.} \]
Check: a skin with the same stress at both edges ($\sigma_2/\sigma_1 = +1$) would give $t_Db/2 = 1500\mm^2$, three
times as much; the web gives less because half of it is in tension and half in compression.

\solutionfor{3.1}{Finish the channel}
Each boom changes the flow by $-10^{-4}\times300\times y_r$:
\begin{align*}
q_{23} &= -6 - 10^{-4}(300)(200) = -6 - 6 = -12\Nmm\\
q_{34} &= -12 - 10^{-4}(300)(-200) = -12 + 6 = -6\Nmm\\
\text{past boom 4} &= -6 - 10^{-4}(300)(-200) = -6 + 6 = 0
\end{align*}
The last one must be zero: the edge beyond boom 4 is free, so it carries no shear. Getting zero there is the check that
every step was right. The minus signs mean the flow runs opposite to the direction $s$.

\solutionfor{4.1}{Practice box: booms, then the shear centre}
(a) Angles: left $(50+50-6)\times6 = 564\mm^2$, right $(40+30-5)\times5 = 325\mm^2$.
Skins (ratio $+1$): $\dfrac{8\times600}{6}(2+1) = 2400\mm^2$ to each corner.
Webs (ratio $-1$): left $\dfrac{6\times240}{6}(1) = 240$, right $\dfrac{5\times240}{6}(1) = 200\mm^2$.
\[ B_{left} = 564 + 2400 + 240 = 3204\mm^2, \qquad B_{right} = 325 + 2400 + 200 = 2925\mm^2 . \]
(b) Symmetric about the horizontal axis, so $I_{xy} = 0$ and
$I_{xx} = 2(3204 + 2925)(120)^2 = 176.5\times10^6\mm^4$.

(c) Booms 1 (top left), 2 (top right), 3 (bottom right), 4 (bottom left); cut wall 12 and go $1\to2\to3\to4\to1$:
\begin{align*}
q_{b,12} &= 0\\
q_{b,23} &= -\frac{S_y}{I_{xx}}(2925)(120) = -1.989\times10^{-3}\,S_y\\
q_{b,34} &= q_{b,23} - \frac{S_y}{I_{xx}}(2925)(-120) = 0\\
q_{b,41} &= 0 - \frac{S_y}{I_{xx}}(3204)(-120) = +2.178\times10^{-3}\,S_y
\end{align*}
(d) At the shear centre the section does not twist: $\oint (q_b + q_{s,0})\,ds/t = 0$, with $ds/t$ = top
$600/8 = 75$, right $240/5 = 48$, bottom $75$, left $240/6 = 40$ (total $238$):
\begin{align*}
&(-1.989\times10^{-3})(48) + (2.178\times10^{-3})(40) + 238\,q_{s,0} = 0\\
&-95.45\times10^{-3} + 87.13\times10^{-3} + 238\,q_{s,0} = 0
\;\Rightarrow\; q_{s,0} = 3.50\times10^{-5}\,S_y .
\end{align*}
Totals: $q_{23} = -1.954\times10^{-3}S_y$ and $q_{41} = 2.213\times10^{-3}S_y$ (both webs flow upward). Web forces:
$1.954\times10^{-3}\times240 = 0.469\,S_y$ and $2.213\times10^{-3}\times240 = 0.531\,S_y$; together $1.000\,S_y$. \checkmark

Moments about the left web: the right web gives $0.469\,S_y\times600 = 281.3\,S_y$; the skins' small flows form an
opposite couple $3.50\times10^{-5}S_y\times600\times240 = 5.0\,S_y$. So
\[ S_y\,\xi = 281.3\,S_y - 5.0\,S_y \;\Rightarrow\; \xi = 276.3\mm \text{ from the left web,} \]
left of the middle ($300\mm$) because the left booms are heavier, as in the homework.

\solutionfor{5.1}{Twist rate of one cell}
Around cell I every wall carries $q_I$ except the shared wall, which carries $q_I - q_{II}$:
\begin{align*}
\frac{d\theta}{dz} &= \frac{1}{2A_IG}\oint q\,\frac{ds}{t}
= \frac{1}{2A_IG}\Big[q_I(\delta_I - \delta_{12}) + (q_I - q_{II})\delta_{12}\Big]\\
&= \frac{1}{2A_IG}\big(q_I\delta_I - q_{II}\delta_{12}\big).
\end{align*}

\solutionfor{6.1}{Practice two-cell box: $T = 15$ kN m}
(a) Webs and spar keep their thickness. Skins get $t^* = (24\,000/27\,000)\,t = 0.8889\,t$, so $0.889\mm$ (cell I)
and $0.711\mm$ (cell II). Then $\delta = L/t^*$:
\begin{align*}
&\text{left web } 250/2.0 = 125, \quad \text{spar } 250/2.5 = 100, \quad \text{right web } 250/1.6 = 156.3,\\
&\text{each cell-I skin } 500/0.889 = 562.5, \quad \text{each cell-II skin } 600/0.711 = 843.8 .
\end{align*}
\[ \delta_I = 125 + 2(562.5) + 100 = 1350, \qquad \delta_{II} = 156.25 + 2(843.75) + 100 = 1943.75 . \]
(b) Areas $A_I = 500\times250 = 125\,000\mm^2$, $A_{II} = 600\times250 = 150\,000\mm^2$. Equal twist:
\begin{align*}
&\frac{1350\,q_I - 100\,q_{II}}{2(125\,000)} = \frac{-100\,q_I + 1943.75\,q_{II}}{2(150\,000)}\\
&\Rightarrow\; 150\,000(1350q_I - 100q_{II}) = 125\,000(-100q_I + 1943.75q_{II})
\end{align*}
\[ \Rightarrow\; 215.0\times10^6\,q_I = 258.0\times10^6\,q_{II} \;\Rightarrow\; q_I = 1.200\,q_{II} . \]
Torque, with $T = 15\times10^6$ N\,mm:
$15\times10^6 = 2(125\,000)(1.200\,q_{II}) + 2(150\,000)\,q_{II} = 600\,000\,q_{II}$, so
\[ q_{II} = 25.0\Nmm, \qquad q_I = 30.0\Nmm, \qquad q_{spar} = q_I - q_{II} = 5.0\Nmm . \]
Both cells flow counter-clockwise, with the torque; in the spar the net flow follows cell I.

(c) $\tau = q/t$ with the real thickness: left web $30.0/2.0 = 15.0$; cell-I skins $30.0/1.0 = 30.0$; spar
$5.0/2.5 = 2.0$; cell-II skins $25.0/0.8 = 31.3$; right web $25.0/1.6 = 15.6\ \mathrm{N/mm^2}$.

(d) \[ G_{REF}\,\frac{d\theta}{dz} = \frac{30.0(1350) - 25.0(100)}{2(125\,000)} = \frac{40\,500 - 2\,500}{250\,000} = 0.152, \]
so $\dfrac{d\theta}{dz} = \dfrac{0.152}{27\,000} = 5.63\times10^{-6}\ \mathrm{rad/mm} = 0.32^\circ/\mathrm{m}$. Check with cell II:
$\dfrac{-30.0(100) + 25.0(1943.75)}{2(150\,000)} = \dfrac{45\,594}{300\,000} = 0.152$. \checkmark
\end{document}
